\section{Introduction}

Contrastive self-supervised learning is widely assumed to amplify spurious correlations. Its augmentation-invariance objective encourages representations where features stable across random crops and color jitter become strongly instance-discriminative --- and salient spurious features like background texture satisfy both criteria. This theoretical account predicts that contrastive SSL backbones should encode spurious attributes \textit{more} strongly than supervised counterparts. We show the opposite is true.

In our experiments on Waterbirds, supervised ERM encodes spurious background features significantly more strongly than contrastive MoCo-v3, with a ratio difference of 2.5 percentage points and an exceptionally large effect size (Cohen's $d = 5.68$, $p < 0.0001$, Bonferroni-corrected). The paradigm ranking is ERM $\approx$ DINO $>$ BarlowTwins $>$ MoCo-v3 --- the ranking an augmentation-invariance account would predict in reverse.

This finding has practical consequences. Practitioners selecting pretraining backbones for fairness-critical or distribution-shift-sensitive applications routinely choose between supervised ImageNet models and self-supervised alternatives without principled guidance on their spurious feature content. If contrastive SSL were the worst offender --- as theory suggests --- the field's move toward SSL for robustness would be counterproductive. If it is the best, that changes the design calculus considerably.

\textbf{The gap this work fills.} Despite substantial prior work on spurious correlations in deep learning \citep{Sagawa2020GroupDRO, Kirichenko2022DFR, Geirhos2020Shortcut}, systematic cross-paradigm comparison of spurious feature encoding in frozen representations is missing. DFR \citep{Kirichenko2022DFR} demonstrates that ERM features are sufficient for robustness after group-balanced retraining, but never compares to SSL. WILDS \citep{Koh2021WILDS} and DomainBed \citep{Gulrajani2021DomainBed} evaluate fine-tuned models across distribution shifts, conflating pretraining representation quality with fine-tuning dynamics. Contrastive shortcut learning work \citep{Robinson2021Contrastive} confirms shortcuts appear in SSL representations, but does not compare their magnitude to supervised baselines with a controlled backbone. To the best of our knowledge, we provide the first controlled 4-paradigm comparison on frozen ResNet-50 features using spurious/task probe accuracy ratio as a scale-free paradigm discriminator.

\textbf{Key insight.} The dominant driver of spurious feature encoding is supervised label correlation, not augmentation-based instance-discriminability. ERM's cross-entropy objective is trained on data where spurious attributes (Waterbirds background: 95\% training correlation) reliably predict class labels, causing the model to jointly encode both task-discriminative features and spurious features. MoCo-v3, being label-agnostic, encodes spurious features only to the degree they are instance-discriminative under the augmentation set --- a weaker pressure that produces measurably lower spurious encoding. DINO's self-distillation, which implicitly generates class-level semantic targets, replicates ERM's feature selection behavior (ERM $\approx$ DINO, $p_\text{bonf} = 1.0$), supporting this interpretation.

\textbf{Contributions.} We make the following contributions:

\begin{enumerate}
\item \textbf{Empirical finding (Waterbirds):} A controlled 4-paradigm $\times$ 5-seed comparison on frozen ResNet-50 demonstrates a highly significant paradigm effect on spurious/task probe accuracy ratio (ANOVA $F = 35.99$, $p = 2.42 \times 10^{-7}$), with ERM encoding spurious background features most strongly and MoCo-v3 least --- opposite to augmentation-invariance theory.

\item \textbf{Empirical finding (cross-dataset):} CelebA replication (ANOVA $F = 5.51$, $p = 0.009$) confirms the paradigm effect across datasets, but reveals a ranking reversal (MoCo-v3 lowest on Waterbirds, highest on CelebA), establishing that spurious attribute type $\times$ pretraining objective interaction is a necessary design consideration.

\item \textbf{Mechanistic insight:} We identify supervised label correlation --- not augmentation invariance --- as the primary driver of spurious feature encoding, supported by the ERM $\approx$ DINO similarity ($d = 0.60$, $p_\text{bonf} = 1.0$) and the ERM $>$ MoCo-v3 directional result ($d = 5.68$, $p = 9.4 \times 10^{-6}$ for reversed direction relative to prediction).

\item \textbf{Preliminary finding:} Background-replacement augmentation in contrastive training is verified as an active causal lever for spurious encoding modulation (mechanism activation confirmed at $19\times$ pixel-difference threshold). Full statistical characterization of the spurious ratio reduction is ongoing; we report this as a proof-of-concept result (see Discussion~\ref{sec:discussion}).
\end{enumerate}

The paper is organized as follows. Section~\ref{sec:related} reviews related work on spurious correlations and representation learning. Section~\ref{sec:method} describes our methodology. Section~\ref{sec:experiments} details the experimental setup. Section~\ref{sec:results} presents results. Section~\ref{sec:discussion} discusses implications and limitations. Section~\ref{sec:conclusion} concludes.
